\cotefile{79}

\scanpage{190}

% B01 / text / top, condition b), lines 1–4
Considérons maintenant la condition b), \deleted{\unsure{Comme}} \inserted{supposons a) satisfaite}. Si $E=\mathrm{Et}_{A_0}(s)$, alors $E$ \deleted{\unsure{est s|est}} stable par $\widetilde U_s$, et l'action de $\widetilde U_s$ sur $E$ transforme $A_{E,\lambda}$ en \unsure{lui-\={m}|lui-même}.
% slanted
\margin{On \unsure{reprendra|reprend|reviendra} le fil de la réflexion d'un nouveau pied, plus bas. --}

% archive sheet number: 189
\scanpage{191}

% B01 / text / top, heading and first line
\mspage{98}\uline{2.}

Soit \deleted{$(S,A_0)$ un graphe}

% B02 / display / upper, (1)
\[
 (S,A_0)\qquad A_0\subset\unsure{P|\mathcal P|\mathfrak P}_2(S) \tag*{(1)}
\]

% B03 / text / upper, before (2)
un graphe, satisfaisant la condition de ``triangularité''

% B04 / display / upper, (2), left brace
\[
\left\{\begin{aligned}
 &\text{a)}\quad S=\bigcup_{a\in A_0}a\\
 &\text{b)}\quad \forall\,\underset{\substack{\shortparallel\\\{s,t\}}}{a}\in A_0,\quad \exists\ \text{exactement deux éléments }u\text{ de }S,\\
 &\hspace{12em}\text{tels que }\{s,u\},\{t,u\}\in A_0
\end{aligned}\right. \tag*{(2)}
\]

% B05 / text / upper middle, before (3)
\deleted{\unsure{avec}} de telle façon que pour tout $s\in S$, son étoile pour $A_0$

% B06 / display / upper middle, (3)
\[
 \mathrm{Et}_{A_0}(s)=\{t\in S\mid\{s,t\}\in A_0\}\qquad =E \tag*{(3)}
\]

% B07 / text / middle, after (3)
est de façon canonique munie d'une structure de contour \deleted{\unsure{dont}} \unsure{indenté|induit|induite}, dont l'ens des arêtes est

% B08 / display / middle, (4)
\[
 A_{0E}=P_2(E)\cap A_0 \tag*{(4)}
\]

% B09 / text / middle, bracketed passage
dont les composantes connexes sont d'ordre $\geq 3$. [Je suppose donné de plus un groupe d'automorphismes
$G_0$ de $(S,A_0)$, transitif \deleted{sur} $\overrightarrow{A_0}$, et \deleted{simplement}\inserted{même} transitif sur l'ens des répères triangulaires
% slanted
\margin{\unsure{descend|devient|demandé} plus bas}

% B10 / display / lower middle, end of bracket, (5)
\[
 \mathrm{Rep}(\unsure{s|S},A_0)\overset{\mathrm{def}}{=}
 \{(s,t,u)\in S^3\mid \{s,t\},\{t,u\},\deleted{\}}\{u,s\}\in A_0\}\,] \tag*{(5)}
\]

% B11 / text / lower middle, restart
\deleted{Je suppose de plus \unsure{satisf.|satisfaite}} Considérons l'application

% B12 / diagram / lower, (6)
\diagram{(6)\qquad\begin{tikzcd}[ampersand replacement=\&,row sep=0pt]
 {\mathrm{Et}_{A_0}:}\& S\arrow[r,dash]\&P(S)\\
 \&s\arrow[r,mapsto]\&\mathrm{Et}_{A_0}(s)
\end{tikzcd}}

% B13 / text / lower, before (7)
et soit $\unsure{\mathcal E|\mathscr E|\Sigma}_A$ son image

% B14 / display / lower, (7), vertical inclusion
\[
 \underset{\substack{\cap\\P(S)}}{\mathcal E_A}=\{E\in P(S)\mid\exists\,s\in S\ \text{t.q.}\ E=\mathrm{Et}_{A_0}(s)\} \tag*{(7)}
\]

% B15 / text / foot, continuing
On note que pour tout $E\in\mathcal E_A$, la structure de contour
sur $E$ définie par (4) ne dépend pas du choix du $s\in S$ tel que $E=\mathrm{Et}_{A_0}(s)$. On récupère $A_0$, à

% archive sheet number: 190
\scanpage{192}

% B01 / text / top, continuation
partir de ces structures sur les $E\deleted{_{\unsure{s}}}\in\unsure{\mathcal E|\mathscr E|\Sigma}$, par la formule

% B02 / display / top, (8)
\[
 A_0=\bigcup_{E\in\mathcal E}A_{0E}\ . \tag*{(8)}
\]

% B03 / text / upper, before (9)
Revenant à l'action de $G_0$ sur $S$, $G_0$ opère
par transport de structure sur les \deleted{objets} ens. introduits, et \deleted{l'}l'application \deleted{\unsure{surject.}} \unsure{surjective|surjection}

% B04 / diagram / upper, (9)
\diagram{(9)\qquad\begin{tikzcd}[ampersand replacement=\&,row sep=0pt]
 S\arrow[r]\&\mathcal E\\s\arrow[r,mapsto]\&\mathrm{Et}_{A_0}(s)
\end{tikzcd}}

% B05 / text / upper middle, following (9)
commute à l'action de $G$. Les ensembles
$S,A_0,\overrightarrow{A_0},\mathrm{Rep}(S,A_0)$, ainsi que

% B06 / display / middle, (10)
\[
 F_0=\{f\in\unsure{P|\mathcal P|\mathfrak P}_3(S)\mid\forall a\in P_2(f),\ a\in A_0\} \tag*{(10)}
\]

% B07 / text / middle, homogeneous spaces
sont des espaces homogènes sous $G$ -- tous sont
des espaces homogènes quotients de $\mathrm{Rep}(S,A_0)$,
qui est un torseur sous $G$.

% B08 / text / lower middle, graph structures
Je m'intéresse maintenant à l'ens des
structures de graphes $(S,A_i)$ sur \unsure{le|la} $\overline{\mathrm m}$ ens. des
sommets $S$, qui admettent $G$ comme groupe
d'automorphismes, sont triangulaires, et donnent
lieu au\deleted{\unsure{x mêmes étoiles que|même ensemble}}\inserted{$\overline{\mathrm m}$ ens. d'étoiles} $\mathcal E$ que $A_0$. \deleted{Soit}

% B09 / text / lower, determining the set
Je voudrais déterminer cet ens \inserted{$\Pi=(\unsure{A'|A^i|A_i})_{i\in\Pi}$} des structures,
qui a priori ne dépend que de $S$, $G_0\subset\unsure{\mathfrak S|\mathfrak G}_S$,
et de $\mathcal E\subset P(S)$ \deleted{\unsure{stable}} fixé par $G_0$, et non
de \deleted{\unsure{la spécif.}} la donnée de $A_0\subset P_2(S)$ -- cette
donnée assure uniquement que $\Pi\ne\varnothing$, puisque

% B10 / display / lower, (11)
\[
 A_0\in\Pi\subset P(P_2(S))\ . \tag*{(11)}
\]

% B11 / text / foot, before (12)
Comme on a dit, \inserted{pour $A_0$, un} $A_i$ est déterminé quand on
connaît les structures induites

% B12 / display / foot, (12)
\[
 A_{iE}=A_i\cap E\qquad(E\in\mathcal E) \tag*{(12)}
\]

% archive sheet number: 191
\scanpage{193}

% B01 / text / top, continuation
\mspage{99}qui sont des structures de contour sur les
$E\in\unsure{\mathcal E|\mathscr E|\Sigma}$, \deleted{\unsure{ces}} structures \unsure{redonnant|redonnent} $A_i$ par

% B02 / display / upper, (13)
\[
 A_i=\bigcup_{E\in\mathcal E}A_{iE}\ . \tag*{(13)}
\]

% B03 / text / upper, before (14)
L'application $\mathrm{Et}_{A_i}$ correspondante

% B04 / diagram / upper, (14)
\diagram{(14)\qquad\begin{tikzcd}[ampersand replacement=\&,row sep=0pt]
 {\mathrm{Et}_{A_i}:}\& S\arrow[r,"\mathrm{surj}"]\&\mathcal E\\
 \&s\arrow[r,mapsto]\&\mathrm{Et}_{A_i}(s)
\end{tikzcd}}

% B05 / text / upper middle, following (14)
dépend a priori de la \deleted{\unsure{donnée|l'indice}} structure choisie $i\in\Pi$.

% B06 / text / middle, transitivity
\deleted{Notons que} Comme $G_0$ opère transitivement sur
$\mathcal E$ et respecte la structure $A_i$, \deleted{$A_i$}\inserted{les $A_{iE}$ \deleted{donc $A_i$}} sont connus,
donc $A_i$ est connu, dès qu'on connait
$A_{iE}$ pour \uline{un} $E_0\in\mathcal E$ :

% B07 / diagram / middle, (15)
\diagram{(15)\qquad L'application \inserted{$\unsure{\psi|\Psi|\varphi}_{E}:$}
\begin{tikzcd}[ampersand replacement=\&,row sep=small]
 \Pi\arrow[r]\&\substack{\text{Structures de contour}\\\text{sur }E_0}\\
 i\arrow[r,mapsto]\& A_{iE_0}=A_i\cap\unsure{P|\mathcal P|\mathfrak P}_2(E_0)
\end{tikzcd}}

% B08 / text / middle below (15)
est injective, et son image est formée
de structures de contour qui sont
\deleted{stables}\inserted{invariées} par l'opération de $G_{0E_0}$ (stabilisateur
de $E_0$ dans $G_0$) sur $E_0$.

% B09 / text / lower, system of contours
Il est clair d'ailleurs que la \deleted{\unsure{ens des}}\inserted{donnée d'une} structures
de contour sur $E_0$, invariée par $G_{E_0}$, est équivalente
à la donnée d'un système \inserted{$(C_E)_E$} de structures
de contour sur les $E\in\mathcal E$, invariées par $G$.

% B10 / text / lower, new paragraph
Il faut maintenant exprimer, en termes
\deleted{de la \unsure{s}} d'une

% B11 / display / foot, (16)
\[
 \begin{aligned}
 &\text{structure de contour }\unsure{C|{\mathbb C}}_{E_0}\text{ sur }E_0,\text{ \unsure{invarie|invariée}}\\
 &\text{par }G_{0E_0},
 \end{aligned} \tag*{(16)}
\]

% B12 / text / foot, last line
la condition qu'elle provient d'un $i\in\Pi$.

% archive sheet number: 192
\scanpage{194}

% B01 / text / top, bullet and cancelled lines
$\bullet$\quad Notons que le fait que $G$ soit transitif sur l'ens. des repères de $(S,A_0)$, signifie aussi que
\deleted{a) $G$ est transitif sur $S$ (condition qui ne dép que pour \unsure{tt|tout} $s\in S$, $G_s$ est}

% B02 / display / upper, (17), brace at left
\[
\left\{\begin{aligned}
 &\text{a) }G\text{ est transitif sur }S\quad\text{(condition indépendante du choix de }A_0\text{)}\\
 &\text{b) }\forall s\in S\quad\text{(ou pour \uline{un} }s_0\in S,\text{ c'est kif kif)}\\
 &\quad\text{l'opération de }G_s\text{ sur }\mathrm{Et}_{A_0}(s)=E\text{ est}\\
 &\quad\text{transitif sur l'ens des repères de}\\
 &\quad\text{la structure de contour }(E,A_{0E}).
\end{aligned}\right. \tag*{(17)}
\]

% B03 / text / middle, polygonality and first marginal note
% slanted
\margin{Cel signifie donc que $A_0$ est associé à une structure de 2-polyèdre régulier triangulaire, et que l'application $\mathrm{Et}_{A_0}$ est injective, sauf \unsure{ds|dans} le seul cas de l'octaèdre où \ldots}
Je vais supposer que cette structure de contour,
\uline{dans le cas de $A_0$}, est connexe i.e. polygonale,
de sorte que cette condition implique que

% B04 / display / middle, (18) and second marginal note
\margin{\unsure{R\kern-.1em b|Rb}}\[
\begin{aligned}
 D_E&=\mathrm{Im}(G_s\longrightarrow\unsure{\mathfrak S|\mathfrak G}_E)\quad\text{est égal au groupe}\\
 &\mathrm{Aut}(E,A_{0E})\text{ tout entier, qui est isomorphe}\\
 &\text{au groupe diédral }D_n,\text{ et contient un}\\
 &\text{unique sous-groupe }D_E^+\ \inserted{=\Lambda_E}\ \text{cyclique d'ordre}\\
 &n.
\end{aligned}\tag*{(18)}
\]
% slanted
\margin{On suppose les sommets $s\in S$ d'ordre $n$ pour $A_0$, ou ce qui revient au $\overline{\mathrm m}$, que les $E\in\unsure{\mathcal E|\mathscr E|\Sigma}$ sont de cardinal $n$ $(n\geq3)$}

% B05 / text / lower, dependence on the choices
A priori, \inserted{par sa définition} ce groupe dépend du choix de $s\in S$, mais
sa description comme $\mathrm{Aut}(E,A_{0E})$ ne dépend
pas \deleted{\unsure{des|du}} choix de $s$ -- elle dépend par contre de
celui de $A_0$. En fait, on voit qu'il n'en
dépend pas, mais seulement des données

% B06 / display / lower, (19) and note
% slanted
\margin{On pourrait faire jouer un rôle \unsure{symétrique à|symétriquement en} $S$ et $\mathcal E$}
\[
 (S,\ \deleted{\unsure{\text{groupe}}}\,G_0\subset\mathfrak S_S,\ \mathcal E\subset\unsure{P|\mathcal P|\mathfrak P}(S)) \tag*{(19)}
\]

% B07 / text / lower, before (20)
Car on a en effet

% B08 / display / lower, (20)
\[
 D_E=\mathrm{Im}(G_E\longrightarrow\mathfrak S_E), \tag*{(20)}
\]

% B09 / text / foot, continuing
(alors que a priori l'image de $G_E\supset G_s$ dans
$\mathfrak S_E$ contient celle de $G_s$ et \deleted{\unsure{pourrait}}\inserted{semblerait}

% archive sheet number: 193
\scanpage{195}

% B01 / text / top, continuation
\mspage{100}pouvoir être différente si $S\xrightarrow{\mathrm{Et}_{A_0}}\unsure{\mathcal E|\mathscr E|\Sigma}$ n'était pas injective). En effet, l'image de $G_E$ étant à priori comprise entre $\mathrm{Aut}(E,A_{0E})$ et $\mathrm{Im}(G_s\mathbin{\unsure{-|\to}}\unsure{\mathfrak S|\mathfrak G}_E)$, donc \unsure{égales|égale} à celles-ci, puisqu'elles sont égales.

% B02 / text / upper, polygon structure
La structure de $n$-polygône \deleted{sur}\inserted{$A_{0E_0}$} $E_0$ est
d'ailleurs entièrement déterminée par la
connaissance d'une paire

% B03 / display / upper middle, (21)
\[
 \underline{\unsure{\omega|w}}_0=\{\omega_0,\omega_0^{-1}\}\subset D_{E_0} \tag*{(21)}
\]

% B04 / text / middle, after (21)
de générateurs de $D_{E_0}$, \deleted{opposés} inverses l'un
de l'autre, engendrant $D_{E_0}^+=\Lambda_{E_0}$ -- c'est donc
une paire de permutations circulaires de
$E_0$, dont la connaissance équivaut à celle de
la structure $n$-polygônale $A_{0E_0}$.

% B05 / text / lower half, crossed block, heading and lemma
\deleted{\unsure{Par la}}\deleted{\unsure{les}}
\deleted{\uline{Lemme}}\quad \deleted{Considérons un ensemble $E_0$ de cardinal
$n$ $(n\in\mathbb N,\ n\geq3)$, muni d'une paire (21) de
permutations circulaires définissant une structure
$n$-polygônale sur $E_0$,} \deleted{Pour tout $\nu\in\mathbb Z/n\mathbb Z$,}
\deleted{et \unsure{l'associant|l'associons} la structure}
\deleted{\unsure{considérons}}\deleted{ de torseur sous}

% B06 / display / lower half, crossed (22), right-hand correspondence
\[
 \deleted{\Lambda_0=\Lambda_{\underline\omega_0}=\mathbb Z/n\mathbb Z\mathbin{\unsure{\wedge|\times}}_{\mu_2}\underline\omega_0}\quad
 \left|\begin{aligned}
 &\deleted{\text{à }\xi\in\Lambda_0,\text{ correspond}}\\
 &\deleted{T_\xi:E_0\xrightarrow{\sim}E_0}
 \end{aligned}\right. \tag*{\deleted{(22)}}
\]

% B07 / text / lower, crossed restart
\deleted{définie par \unsure{$A$}} \deleted{Pour tout \unsure{ent} $\nu\in\mathbb Z/n\mathbb Z$
considérons l'homomorphisme de groupes}

% B08 / display / lower, crossed (23)
\[
 \deleted{\alpha_\nu:\xi\longmapsto\nu\cdot\xi^{\unsure{*|\vphantom{1}|\nu}}:\Lambda_0\longrightarrow\Lambda_0} \tag*{\deleted{\unsure{(23)|(22)}}}
\]

% B09 / text / foot, crossed last line and insertion
\deleted{et l'opération \inserted{$T^{(\nu)}$} de $\Lambda_0$ sur $E_0$, définie par}

% B10 / display / foot, crossed formula
\[
 \deleted{T^{(\nu)}_{\xi}(x_0)=T^{\unsure{\nu|(\nu)}}_{\deleted{\unsure{\xi}}\alpha_\nu(\xi)}(x_0)=(T_\xi)^\nu(x)}
\]

% archive sheet number: 194
\scanpage{196}

% B01 / text / top, proposition and additions over first lines
\deleted{Considérons}
\uline{Proposition}\quad Soit $E$ un torseur sous un
groupe $\Lambda_E$ \inserted{$\subset\unsure{\mathfrak S|\mathfrak G}_E$} cyclique d'ordre $n$. \inserted{$(n\in\mathbb N,\ \deleted{\unsure{3|2}}\ n\geq\unsure{2|3})$} Considérons
l'ens. $\overrightarrow C(E)$ des structures de contour
\deleted{condition}\inserted{\uline{orienté}} sur $E$, invariées par $\Lambda_E$.

% B02 / text / upper middle, bijection and first marginal note
\deleted{On a} \deleted{Pour tout} \deleted{Pour tout}
Cet ens est en correspondance biunivoque
avec $\Lambda_E$, en associant à tout $u\in\Lambda_E$
l'unique structure de contour orienté, pour
laquelle $u$ soit l'opération de ``rotation \inserted{génératrice}'' associée à
l'orientation.
% slanted
\margin{$\Lambda_E$ est simplement transitif sur l'ens des repères directs de chacune des structures envisagées (à l'exception des deux cas \unsure{$\mathrm{tot}^{\mathrm t}$|tout à fait} dégénérés, en décomposant en 1-gônes ou bigônes --)}

% B03 / text / middle, proof and NB in the left margin
En effet, la donnée d'une structure $\overrightarrow C$
de contour orienté, sur un ens fini \inserted{$E$}, revient
à celle d'un $u\in\mathfrak S_E$ -- la rotation \unsure{générique|génératrice} du contour orienté. Si un groupe
$\Lambda_E$ opère sur $E$, dire que $\overrightarrow C$ est invariante
par ce groupe, signifie que \inserted{$u=u_{\overrightarrow C}$} l'est, i.e.
que $u$ commute à ce groupe. Or dans le
cas d'une structure de torseur, le
commutant est réduit à $\Lambda_E$ -- cqfd.
% slanted
\margin{\uline{NB} Ici on n'exclut pas les contours ``\unsure{ayant}'' des orbites d'ordre 1 et 2 -- ce qui correspond à une définition \unsure{intuitible|intuitive|intraitable|inutilement restrictive} dans l'optique où la structure est portée par l'ens des sommets \ldots Si on voulait l'éviter, il faudrait dans la proposition se borner aux $u\in\Lambda_E$ qui ne sont pas d'ordre 1 ou 2, ce qui revient à \inserted{éliminer} juste les \deleted{\unsure{sommets}} éléments du \unsure{ss-groupe|s-groupe} d'ordre 2 \inserted{$ {}_2\Lambda_E$}}

% B04 / text / foot, corollary 1
\uline{Corollaire 1} Les structures $C$ envisagées forment
un module libre de rg 1 sur $\Lambda_n=\mathbb Z/n\mathbb Z$
(savoir $\Lambda_E$ !). Les générateurs de ce
module correspondent aux structures
polygônales orientées sur $E$, invariées par $\Lambda_E$,
qui forment donc un \uline{torseur} \inserted{$\Lambda_E^*$} sur $(\mathbb Z/n\mathbb Z)^*$.

% archive sheet number: 195
\scanpage{197}

% B01 / text / top, corollary 2 and NB
\mspage{101}\uline{Corollaire 2}\quad Toute structure \deleted{orientée} de contour $C$
sur $E$, \unsure{invarie|invariée} par $\Lambda_E$, admet exactement
deux orientations invariées par $\Lambda_E$, elles sont
opposées l'une à l'autre. De cette façon, l'ens
$\unsure{\mathcal C|\mathscr C}(E)$ de ces contours s'identifie à l'ensemble
quotient
% slanted
\margin{\uline{NB} Le sous groupe \deleted{\unsure{$D_E$ de}} $\unsure{\mathfrak S|\mathfrak G}_E$ qui normalise $\Lambda_E$ est le ss-groupe affine $\mathrm{Aff}_{\Lambda_E}(E)$, et $D_E^+$ se décrit comme l'image inverse de $\{\pm1\}\subset(\mathbb Z/n\mathbb Z)^*$ dans cette ext. de $(\mathbb Z/n\mathbb Z)^*$ par $\Lambda_E$. Les structures de contour envisagées invariées par $\Lambda_E$ sont alors $\overline{\mathrm m}$ invariées par $D_E$ (mais non par $\mathrm{Aff}(E)$ tout entier, bien sûr). Le groupe $D_E$ est alors simplement transitif sur l'ens des répères de toute structure \deleted{\unsure{polyg}} de contour envisagée (à l'exclusion des deux dégénérées)}

% B02 / display / upper middle, (22)
\[
 \mathcal C(E)\simeq\bigl(\Lambda_E\smallsetminus\{0,\overset{\inserted{\text{si }n\text{ pair}}}{\sigma}\}\bigr)/\mathord{\pm}\mathrm{id}_{\Lambda_E} \tag*{(22)}
\]

% B03 / text / middle, parenthesis
(\deleted{On} \deleted{prend} l'ensemble $\Lambda_E/\mathord{\pm}\mathrm{id}_{\Lambda_E}$ plus sympathique,
si on admet les deux contours invariés ``impropres'',
correspondants à $0$ et $\sigma$ $\ldots$) dans ce cas,

% B04 / text / middle, monoid and cancellations
\deleted{\unsure{$\Lambda(\mathbb Z)$}}\deleted{le groupe} \inserted{le monoïde} $\Lambda_n^\times=\mathbb Z/n\mathbb Z$ avec la multiplication
opère sur cet ensemble, qui \deleted{\unsure{doit admet une|admet un}}\inserted{est \deleted{\unsure{ainsi}}} sera
\deleted{pour} \unsure{non can.$^{\mathrm t}$|non can.|non com.$^{\mathrm t}$} \unsure{à|;} $\Lambda_n^\times$ comme ens à opérateurs \unsure{?}

% B05 / text / lower, polygonal structures
L'ens des structures \deleted{de} polygônales sur $E$
invariées par $E$ \deleted{forme un torseur}\inserted{est en corresp 1--1} avec
$\Lambda_E^*/\mathord{\pm}1$, et forme un torseur sous \deleted{\unsure{$\Lambda'_n=\Lambda$}}

% B06 / display / lower, (23)
\[
 \unsure{Z'_n|\mathbb Z'_n}=(\mathbb Z/n\mathbb Z)^*/\mathord{\pm}1\ . \tag*{(23)}
\]

% B07 / text / foot, operators
Pour $\unsure{\nu|v}\in(\mathbb Z/n\mathbb Z)^*$ resp $\unsure{\dot\nu|\dot v}\in Z'_n=(\mathbb Z/n\mathbb Z)^{\unsure{*|1}}/\mathord{\pm}1$,
je désigne par $T_\nu$ resp $T_{\dot\nu}$ l'opération
correspondante sur l'ens $\overrightarrow{\mathcal C}_E$ resp. $\widehat{\mathcal C}_E$ --
et cette \unsure{dénotation|denotation|opération} \unsure{les les identiant|les les identifiant|des les identiant|les éléments} pourra
être prolongée au cas de $\nu\in\mathbb Z/n\mathbb Z^{\unsure{*|\vphantom{1}}}$ (sans $*$ !)
$\dot\nu\in(\mathbb Z/n\mathbb Z)/\mathord{\pm}1$ --

% archive sheet number: 196
\scanpage{198}

% B01 / text / top, powers of a generator
En termes d'un générateur $\unsure{\omega|w}_0$ de $\Lambda_E$,
i.e. $\omega_0\in\Lambda_E^*$, on trouve la structure de
contour orienté correspondant à $\unsure{\nu|v}\cdot[\omega_0]=[\omega_0^\nu]$,
via $u=\omega_0^\nu$ -- et pour $\dot\nu=\{\pm\nu\}$, on
prend la structure de contour définie par la
paire $\{\omega_0^\nu,\omega_0^{-\nu}\}$.

% B02 / text / upper middle, return to the situation
Revenons maintenant au cas de la situation
$(S,G_0,\subset\unsure{\mathfrak S|\mathfrak G}_S,\ \unsure{\mathcal E|\mathscr E|\Sigma}\subset\unsure{P|\mathcal P|\mathfrak P}_E)$, et on voit donc que les
systèmes
% slanted
\margin{\uline{NB} On voit maintenant que $G$ est transitif sur l'ens des}

% B03 / display / middle, system of contours
\[
 (C_E)_{E\in\mathcal E}\qquad
 \begin{aligned}
 &C_E\in\text{ structures de}\\
 &\hspace{2em}\text{contour sur }E\\
 &C_E=(E,A_E)
 \end{aligned}
\]

% B04 / text / middle, parametrization, left margin continuation
invariés par $G$, correspondent biunivoquement
aux éléments \inserted{\deleted{\unsure{$u$}}$\unsure{v|\nu}_{E_0}$} de
$\bigl(\Lambda_{E_0}\smallsetminus\underbrace{\{0,\sigma\}}_{\text{sous-groupe d'ordre 2}}\bigr)/\mathord{\pm}1$.
% slanted
\margin{ou \unsure{au|aux|on} syst. d'éléments $(v_E)$ \unsure{aux|avec} $v_E\in(\Lambda_E\smallsetminus\{0,\sigma_E\})/\mathord{\pm}1$ \unsure{associés|invariés} par $G\ldots$}
\deleted{\unsure{En}}\deleted{Nous} \unsure{$\ldots$}

% B05 / text / lower middle, realization question
Nous proposons maintenant, \deleted{pour} un tel
$v_{E_0}$ ou $v_E$, d'examiner s'il provient bien
d'une \inserted{structure} $i\in\Pi$.

% B06 / text / lower, assigning the edge set
\deleted{Supposons donc}
Considérons donc le système \deleted{$(C_E)$}
$C_E=(E,A_E)$ associé à un $v_{E_0}$ (ou à un $(v_E)_{E\in\mathcal E}$)
et posons

% B07 / display / lower, (24)
\[
 A=\bigcup_{E\in\mathcal E} A_E\qquad\subset P_2(\unsure{S|B}), \tag*{(24)}
\]

% B08 / text / foot, transitivity
il est clair que $A$ est stable sous $G$,
et que $G$ est transitif sur l'ens des
arêtes orientées de la structure de graphe $(S,A)$.

% archive sheet number: 197
\scanpage{199}

% B01 / text / top, section 3
\mspage{102}3.\quad Je vais reprendre la notion \inserted{de} graphe et de graphe triangulaire,
dans une optique un peu différente. Une structure
de graphe strict peut être interprétée comme la donnée
d'un \deleted{\unsure{gr}} ens. de sommets $S$ et d'une partie
\uline{symétrique} de $S\times S\smallsetminus S$

% B02 / display / upper, (1) and NB to the left
\[
 R_0\subset S\times S\unsure{\smallsetminus|\vphantom{S}},\quad R_0\cap\mathrm{diag}=\varnothing,\quad{}^sR=R, \tag*{(1)}
\]
% slanted
\margin{\uline{NB} si on n'exige pas $R_0\cap\mathrm{diag}=\varnothing$, on trouve une notion un peu plus générale, avec $A_0\subset\unsure{P|\mathcal P|\mathfrak P}_{\mathrm{card}\in\{1,2\}}(S)$, qui correspond à celle d'un graphe pouvant avoir des boucles fermées (une au plus par sommet)}

% B03 / text / upper, before (2)
l'étoile d'un $s\in S$ s'écrivant alors comme

% B04 / display / upper, (2), notation question on right
\[
 \mathrm{Et}_{A_0}(s)=R_0(s)\qquad \text{on note aussi }\mathrm{Et}_{R_0}(s)\ ? \tag*{(2)}
\]

% B05 / text / upper, before (3)
\unsure{puis|que|pour} l'on a

% B06 / display / upper middle, (3)
\[
 A_0=\{\{\unsure{s,t|s;t}\}\in P_2(S)\mid(s,t)\in R_0\}\ . \tag*{(3)}
\]

% B07 / text / middle, restart before (4)
\deleted{Dans le cas où} \inserted{Considérons} l'application

% B08 / diagram / middle, (4)
\diagram{(4)\qquad\begin{tikzcd}[ampersand replacement=\&,row sep=0pt]
 S\arrow[r,mapsto]\&P(S)\\ s\arrow[r,mapsto]\&R_0(s)=\mathrm{Et}_{A_0}(s)
\end{tikzcd}}

% B09 / text / middle, introduction of the image
\deleted{est injective,} désignons \deleted{en tous cas} par $S^*$ son \deleted{\unsure{image}}\inserted{image},
(notée $\unsure{\mathcal E|\mathscr E|\Sigma}$ précédemment

% B10 / display / middle, (5)
\[
 S^*=\{E\in P(S)\mid\exists\,s\in S,\ \text{t.q. }E=R_0(s)\}, \tag*{(5)}
\]

% B11 / text / middle, before (6)
on note que l'inclusion identique

% B12 / display / middle, (6), inclusion and assignment
\[
 S^*\lhook\joinrel\longrightarrow P(S)\qquad s^*\longmapsto E_{s^*} \tag*{(6)}
\]

% B13 / text / lower middle, interpretation
s'interprète aussi comme définissant une relation

% B14 / display / lower middle, (7)
\[
 \deleted{\unsure{\rho|P}}\unsure{R_*|R_{\star}|\vec{R}_*}\subset S\times S^* \tag*{(7)}
\]

% B15 / text / lower, before (8)
liée à (6) par

% B16 / display / lower, (8)
\[
 \forall\,s^*\in S^*,\quad\text{on a}\quad{}^sR_*(s^*)=E_{s^*}\ . \tag*{(8)}
\]

% B17 / text / lower, bracketed special case
[\deleted{d'où} Nous nous intéressons surtout au cas où

% B18 / display / lower, crossed (9)
\[
 \deleted{\forall\,s\in S,\ \text{on a }R_0(s)} \tag*{\deleted{(9)}}
\]

% B19 / text / lower, (4) injective
(4) est injective (\deleted{\unsure{tout}}\inserted{\unsure{une|un} ($A_0$-)} étoile détermine son centre),
et où on \deleted{peut écrire}\inserted{trouve} alors, pour tout $s\in S^*$, une
application bijection canonique

% B20 / diagram / foot, (9)
\diagram{(9)\qquad\begin{tikzcd}[ampersand replacement=\&,row sep=small]
 \underbrace{R_*(s)}_{\substack{\text{ens des étoiles}\\ E\ \text{qui contiennent }s}}\&\mathrm{Et}_{A_0}(s)\arrow[l,"\sim"']\\
 \mathrm{Et}_{A_0}(t)\&t\arrow[l,mapsto]
\end{tikzcd}}

% B21 / text / foot, last parenthesis
(en général, on a seulement une application surjective)]

% archive sheet number: 198
\scanpage{200}

% B01 / text / top, restart
\deleted{Posons}\inserted{Écrivons} maintenant indiciellement l'application (4) :

% B02 / display / top, (10)
\[
 \mathrm{Et}_{A_0}(s)=E_{\varphi_{A_0}(s)}={}^{s}_{\unsure{0|\vphantom{1}}}R_*(\varphi(s)) \tag*{(10)}
\]

% B03 / text / upper, before (11)
où

% B04 / display / upper, (11)
\[
 \varphi_{\unsure{0A_0|0|A_0}}:S\xrightarrow{\unsure{\sim|\mathrm{su}}}S^* \tag*{(11)}
\]

% B05 / text / upper, bracketing and restart
est l'application \deleted{\unsure{b}} surjective [en fait bijective
\deleted{par}\inserted{vu} l'hypothèse \deleted{faite}\inserted{envisagée plus haut}) \deleted{de} induite par (4), \deleted{\unsure{on}}
$[\![$\deleted{justement définie}
\inserted{Ainsi la formule (9)} s'écrit aussi

% B06 / display / upper middle, (12)
\[
 \forall s\in S,\quad \underset{\substack{\cap\\S^*}}{R_*(s)}\xleftarrow{\sim}E_{\varphi_0(s)}\ .\quad ]\!] \tag*{(12)}
\]

% B07 / text / middle, injectivity
Ainsi, la structure \inserted{de graphe strict} $(S,A)$, $A\subset\unsure{P|\mathcal P|\mathfrak P}_2(S)$, $[\![$satisfaisant
la condition d'injectivité

% B08 / display / middle, (13), left brace and marginal NB
\[
 \left\{\begin{aligned}
 &s\longmapsto\mathrm{Et}_{A_0}(s)\quad\text{de } S\text{ dans }P(S)\text{ injective}\\
 &\text{i.e.}\quad \forall\,s\!,t\in S,\ s\ne t,\quad\exists\,u\in S\text{ tel que l'on ait}\\
 &\qquad(\{s,u\}\in A_0,\{t,u\}\notin A_0)\quad\text{ou l'inverse}\\
 &\qquad(\{t,u\}\in A_0,\{s,u\}\notin A_0)\\
 &\text{i.e.}\quad \forall\,a\in P_2(S),\ \exists\,u\in S,\ s,t\in a,\text{ tels que}\\
 &\qquad\{s,u\}\in A_0,\{t,u\}\notin \unsure{a_0|A_0}\quad ]\!]
 \end{aligned}\right. \tag*{(13)}
\]
% slanted
\margin{\uline{NB} Même sans condition d'injectivité (13), la structure de graphe $(S,A_0)$ définit une structure (14), à cela près que $\varphi_0$ n'est plus néc. bijective, mais seulement surjective. Mais pour la suite cela se trouve être rédhibitoire -- brisant la symétrie entre $S$ et $S'$, en nous empêchant de passer à $\varphi_0^{-1}$ \ldots C'est dès qu'on veut dualiser qu'on est obligé de faire l'hypothèse \unsure{séparants|séparante} \ldots}

% B09 / text / middle, before (14)
définit les structures suivantes

% B10 / display / lower middle, (14), left brace
\[
 \left\{\begin{aligned}
 &\text{a) une partie } S^*\subset P(S),\quad\text{donc un \deleted{\unsure{gr}} ens. }S^*\\
 &\quad\text{et une relation}\\
 &\qquad R_*\subset S\times S^*\\
 &\text{b) une application }[\![\text{bijection}]\!]\\
 &\qquad\varphi_0:S\xrightarrow{\unsure{\sim|\mathrm{su}}}S^*,
 \end{aligned}\right. \tag*{(14)}
\]

% B11 / text / lower, reconstruction
permettant de reconstruire \inserted{la relation \deleted{\unsure{$R$}} $R_0\subset S\times S$, et donc la}
donnée structure de graphe $(S,A_0)$, par la formule

% B12 / display / lower, (15), with gloss underneath
\[
 \forall\,s,t\in S,\quad
 \underset{\substack{[\text{i.e. }s\in\underbrace{\mathrm{Et}_{A_0}(t)}_{={}^s R_0(\varphi_0(t))}]}}{(s,t)\in R_0}
 \Longleftrightarrow \deleted{\unsure{t|s}}(s,\varphi_0(t))\in R_* \tag*{(15)}
\]

% B13 / text / lower, before (15 prime)
i.e.

% B14 / display / lower, (15 prime)
\[
 R_0=(\mathrm{id}_S\times\varphi_0)^{-1}(R_*), \tag*{$(15')$}
\]

% B15 / text / foot, before (16)
d'où

% B16 / display / foot, (16)
\[
 A_0=\left\{a\in P_2(S)\mid\exists\,\unsure{\{s,t\}|s,t|(s,t)}\in\deleted{\unsure{R_0}}S,\ \text{avec }\left\{\begin{aligned}a&=\{s,t\}\\(s,\varphi_0(t))&\in R_*\end{aligned}\right.\right\}. \tag*{(16)}
\]

% archive sheet number: 199
\scanpage{201}

% B01 / text / top, structure (14)
\mspage{103}$\cdot$\quad Pour que la structure (14)

% B02 / display / top, (17), glosses on right
\[
 (S,S^*,R_*,\varphi_0)\quad\left|\begin{aligned}
 &S\ \text{``ens des sommets''}\\
 &S^*\ \text{ens des ``étoiles'' ou ``cosommets''}\\
 &R_*\ \text{relation ``d'incidence''}\\
 &\varphi_0\ \text{\deleted{application} }\underset{\inserted{\text{jauge}}}{\text{application}}
 \end{aligned}\right. \tag*{(17)}
\]

% B03 / text / upper, conditions for realization
provienne, à isom près, d'une structure de graphe \inserted{$A_0\subset\unsure{P|\mathcal P|\mathfrak P}_2(S)$}
sur $S$, \deleted{\unsure{ou d'un|i.e. d'une}} satisfaisant (13), il faut évidemment
qu'elle satisfasse les conditions suivantes

% B04 / display / upper middle, (18), brace covering a and b
\[
 \left\{\begin{aligned}
 &\text{a) Les applications}\\
 &\qquad s^*\longmapsto{}^sR_*(s^*)\qquad S^*\longrightarrow P(S)\\
 &\qquad s\longmapsto R_*(s)\qquad S\longrightarrow P(S^*)\\
 &\qquad\text{sont \uline{injectives}}\\
 &\text{b) L'application }\varphi_0:S\longrightarrow S^*\text{ est \uline{bijective}, et}\\
 &\quad\text{satisfait la condition de symétrie}
 \end{aligned}\right. \tag*{(18)}
\]
% slanted
\margin{\uline{NB} Si on n'exige pas que $(S,A_0)$ satisfasse la condition (13), il faut laisser tomber dans a) l'injectivité de $S\longrightarrow P(S^*)$, et dans b) l'injectivité de $\varphi_0$ -- $\varphi_0$ sera seulement surjective.}

% B05 / text / middle, symmetry and NB
La relation $(s,\varphi_0(t))\in R_*$ en $(s,t)$ est symétrique
i.e. $\foralls,t\in S$, on a \deleted{$s\ne t$} l'équivalence
\deleted{\unsure{(19)|(18)}}\margin{\uline{NB} On peut remplacer la donnée de $\varphi_0$, plus symétriquement, par la donnée d'une 2\textsuperscript{ième} relation $\Phi_0\subset S\times S^*$, bijective $\ldots$}

% B06 / display / middle, (19)
\[
 (s,\varphi_0(t))\in R_*\Longleftrightarrow(t,\varphi_0(s))\in R_* \tag*{(19)}
\]

% B07 / text / middle, c), and parenthesis
c) \deleted{La relation envisagée \unsure{en}}\inserted{\unsure{est sans}} $\forall s\in S$, on a $(s,\varphi_0(s))\notin R$.
(L'injectivité de $\varphi_0:S\longrightarrow S^*$ i.e. de l'application composée
$S\longrightarrow S^*\longrightarrow P(S)$ n'est autre que l'hypothèse d'injectivité (13).

% B08 / text / middle below, reciprocal
La réciproque est évidente. Ainsi, on trouve
\deleted{Nous avons disti}

% B09 / diagram / lower, (20), groupoid equivalence
\diagram{(20)\quad Équivalence de groupoïdes\par
\begin{tikzcd}[ampersand replacement=\&,column sep=large]
 {\substack{\text{(A)}\\\text{catégorie-groupoïde}\\\text{des graphes stricts}\\(S,A_0)\text{ satisfaisant (13)}\\\text{(sommets déterminés par}\\\text{leurs étoiles)}}}\arrow[r,"\approx"] \&
 {\substack{\text{(B)}\\\text{catégorie-groupoïdes des}\\\text{systèmes }\deleted{\unsure{(S,R)}}(S,S^*,R_*)\\\text{``sommets-cosommets-incidence''}\\\text{\deleted{satisfaisant} condition}\\\text{d'inj. symétrique, munie}\\\text{d'une }\varphi_0\,\text{application bijection}\\S\xrightarrow{\varphi_0}S^*\quad\text{``symétrique''}\\\text{\deleted{et} ``antiréflexive'' }((s,\varphi_0(s))\notin R_0)}}
\end{tikzcd}}

% B10 / text / lower, duality and left-hand gloss
La description (B) fait apparaître une \uline{dualité} dans
la catégorie des graphes envisagés, (``graphes stricts séparés\deleted{\unsure{ants}}''), via
% slanted
\margin{``séparé'' signifie que la notion d'étoile d'un sommet permet de ``séparer'' les sommets''}

% B11 / display / lower, (21), terms renamed below
\[
 (S,S^*,R_*,\varphi_0)\longmapsto
 (\underset{\substack{\shortparallel\\S'}}{S^*},
 \underset{\substack{\shortparallel\\S'{}^*}}{S},
 \underset{\substack{\shortparallel\\R'_*}}{{}^sR_*},
 \underset{\substack{\shortparallel\\\varphi'_0}}{\varphi_0^{-1}}), \tag*{(21)}
\]

% B12 / text / foot, dual graph
d'où la notion de ``graphe dual'' d'un graphe
strict séparé ; mais celui-ci est \uline{canoniquement} isomorphe
au graphe de départ, par

% B13 / display / foot, (22) and right marginal note
\[
 S\xrightarrow{\varphi_0}_{\sim} S^*\qquad S^*\xrightarrow{\varphi_0^{-1}}_{\sim}S \tag*{(22)}
\]
\margin{\uline{NB} La condition de \uline{symétrie} (18) b) signifie justement que (22) est compatible avec les relations $\ldots$}

% archive sheet number: 200
\scanpage{202}

% B01 / text / top, previous description
La description précédente de la notion de graphe a
l'avantage principal d'être ``autoduale'', \uline{et} de
séparer deux composantes de la structure :

% B02 / display / upper, (22), parts a and b
\[
 \begin{aligned}
 &\text{a) La }\unsure{\text{données}|\text{donnée}}\quad(S,S^*,R_*\subset S\times S^*)\\
 &\text{b) La donnée}\qquad \varphi_0:S\xrightarrow{\sim}S^*\ .
 \end{aligned}\tag*{(22)}
\]

% B03 / text / upper middle, without separation
Rien n'empêche d'ailleurs d'étudier de telles
structures, sans supposer que la relation a)
soit ``biséparante'' -- i.e. en ne la soumettant
à aucune hypothèse particulière. La \inserted{auto}dualité
(21) reste valable.

% B04 / text / middle, new question
Le jeu nouveau, c'est qu'on peut partir d'une structure a)

% B05 / display / middle, (23)
\[
 (S,S^*,R_*\subset S\times S^*), \tag*{(23)}
\]

% B06 / text / middle to lower, pregraph
sans donnée d'un $\varphi_0$, -- donc une structure
en un sens plus faible qu'une structure de
graphe, on pourrait l'appeler structure de
``pré-graphe'' -- structure considérée comme ``préliminaire''
à celle de graphe (séparé strict). Il se pose alors
la question d'étudier toutes les structures
de graphe correspondantes i.e. \uline{toutes} les
applications \unsure{bijections|bijectives} \deleted{(ou \unsure{$R_0$|$\mathcal R_0$}}

% B07 / display / lower, (24), left brace
\[
 \left\{\begin{aligned}
 &\varphi:S\xrightarrow{\sim}S^*\\
 &\text{telles que }\varphi\text{ soit ``}R_*\text{-\unsure{symétriques|symétrique} et antiréflexive''}\\
 &\quad\text{(on dit que }\varphi\text{ est une }R_*\text{-jauge)}
 \end{aligned}\right.\quad ; \tag*{(24)}
\]
% slanted
\margin{antiréflexivité n'est pas essentielle, si on admet des graphes avec boucles}

% B08 / text / foot, gauges and origin
\deleted{i.e. les}\inserted{l'ens des} ``$R_*$-jauges''. On s'intéresse surtout au cas,
bien sûr, où cet ensemble n'est pas vide -- et le $\varphi_0$
de tantôt jouera le rôle d'une ``origine'' dans l'ens

% archive sheet number: 201
\scanpage{203}

% B01 / text / upper edge
\mspage{104}
\deleted{\unsure{$\Pi$|$\varphi$}}

% B02 / display / top, (25)
\[
 \Pi=\Pi_R=\{\,\varphi:S\simeq S^*\mid\varphi\text{ est }R\text{-symétrique}\,\} \tag*{(25)}
\]

% B03 / text / upper, choice of gauge
des $R$-jauges. \deleted{On a \inserted{dans} dans ceci.} On voit que
( pour un tel choix $\varphi$

% B04 / display / upper, (26)
\[
 \varphi_0:S\simeq S^*\qquad\varphi_0\in\Pi\ , \tag*{(26)}
\]

% B05 / text / upper to middle, bijections
les applications bijectives $\varphi:S\simeq S^*$ correspondent
biunivoquement aux permutations $\lambda\in\unsure{\mathfrak S|\mathfrak G}_{S}$,
via

% B06 / display / upper middle, (27), permutation below lambda
\[
 \begin{array}{c@{\quad}l}\lambda &\longmapsto\ \varphi=\varphi_\lambda=\varphi_0\circ\lambda\\[-2pt]
 \rotatebox{-90}{$\in$}&\\[-4pt]\mathfrak S_S&\end{array} \tag*{(27)}
\]

% B07 / text / middle, second permutation
ou encore aux $\lambda^*\in\mathfrak S_S$

% B08 / display / middle, (27 star)
\[
 \begin{array}{c@{\quad}l}\lambda^*&\unsure{\longmapsto|\vdash}\ \varphi=\varphi_{\lambda^*}=\deleted{\varphi}\lambda^*\circ\varphi_0\\[-2pt]
 \rotatebox{-90}{$\in$}&\\[-4pt]\mathfrak S_{S^*}&\end{array} \tag*{$(27^*)$}
\]

% B09 / text / middle, introducing (28)
la relation entre $\lambda$ et $\lambda^*$ étant donnée par

% B10 / display / middle, (28), brace
\[
 \left\{\begin{aligned}
 \lambda^*&=\varphi_0(\lambda)=\varphi_0\lambda\varphi_0^{-1}\\
 \lambda&=\varphi_0^{-1}(\lambda^*)=\varphi_0^{-1}\lambda^*\varphi_0
 \end{aligned}\right. \tag*{(28)}
\]

% B11 / text / lower middle, condition
Ceci dit, si $\varphi$ correspond à $\lambda$, pour que
$\varphi$ soit une $R$-jauge, il f. et s. que l'on ait

% B12 / display / lower, (29), brace and interlinear formula
\[
 \left\{\begin{aligned}
 &\text{la relation}\quad(s,\overset{\inserted{\lambda^*\varphi_0(t)}}{\overbrace{\varphi_0(\lambda(t))}})\in R
     \quad\text{est symétrique en }s,t,\\
 &\text{i.e. la relation }\{s,\lambda(t)\}\in A_0\quad\text{symétrique en }(s,t)\\
 &\deleted{\text{est}}\ \text{elle est antiréflexive i.e.}\quad
       \forall s\in S,\ \text{on a }\{s,\lambda(s)\}\notin A_0.
 \end{aligned}\right. \tag*{(29)}
\]
% slanted
\margin{antiréflexivité pas essentielle, si on \unsure{accepte|accept.|exceptées} des graphes impropres}

% B13 / text / lower, symmetry
La symétrie s'écrit aussi sous la forme

% B14 / display / lower, (29 prime), equivalence and implication
\[
 \begin{aligned}
 &\forall(s,t)\in S\times S,\ \text{on a }\{s,t\}\in A_0
 \Longleftrightarrow\{\lambda^{-1}s,\lambda^{\deleted{\unsure{-1|1}}}t\}\in A_0\\
 &\quad(\text{il suffit que }\forall s,t\in S\times S,\quad
 \{s,t\}\in A_0\Longrightarrow\{\lambda^{-1}s,\lambda t\}\in A_0 .
 \end{aligned}\tag*{$(29')$}
\]

% B15 / text / lower left, before last display
\deleted{ou encore} ou \deleted{\unsure{encore|en}}

% B16 / display / foot, (29 double prime)
\[
 \forall s\in S,\ \text{on a }\quad
 \lambda\underbrace{(\mathrm{Et}_{A_0}(\unsure{s|t}))}_{\varphi_0(s)}
 =\underbrace{\mathrm{Et}(\lambda^{-1}s)}_{\varphi_0(\lambda^{-1}s)} \tag*{$(29'')$}
\]

% B17 / text / bottom, unfinished sentence
-- sous cette forme, est bien mis en évidence l'action

% archive sheet number: 202
\scanpage{204}

% B01 / text / top, continuation and struck alternative
de $\lambda$ sur les étoiles \inserted{$E=R_{s^*}$} \deleted{-- elle se fait comme de juste par transport de structure}

% B02 / display / top, struck label and display
\[
 \deleted{\lambda(E)=\lambda(R_{s^*})=}\tag*{\deleted{(30)}}
\]

% B03 / text / upper, new observation
Notons que \deleted{\unsure{$\varphi$|si}} le couple

% B04 / display / upper, (30)
\[
 (\lambda,\lambda^*)\in\unsure{\mathfrak S|\mathfrak G}_{S}\times\mathfrak S_{S^*} \tag*{(30)}
\]

% B05 / text / upper, introduction of condition
défini par \deleted{\unsure{les|de|deux}} couple $(\varphi_0,\varphi)$ de \unsure{g.|q.}$R$-gauges,
satisfait la condition

% B06 / display / upper middle, (31), struck intermediate forms
\[
 \begin{aligned}
 &\deleted{(s,s^*)\in R\Longleftrightarrow(s,\unsure{\varphi_0^{-1}(s^*)|\varphi_0(s^*)})\in R_0
 \Longleftrightarrow\lambda^{-1}s,\lambda\varphi}\ ,\\
 &\deleted{(\lambda s,(\lambda^*)^{-1}s^*)\in R}\\
 &(s,s^\vee)\in R\Longleftrightarrow(\lambda^{-1}s,\lambda^\vee s^\vee)\in R .
 \end{aligned}\tag*{(31)}
\]

% B07 / text / upper middle, explanation
\deleted{\unsure{Soit|et soit}} qui \deleted{ne fait} (pour $\varphi_0$ fixée) ne fait qu'exprimer la condition
$(29')$ sur le $\varphi=\lambda^\vee\varphi_0=\varphi_0\lambda:S\longrightarrow S^\vee$ définie par \deleted{\unsure{$\varphi_0\varphi$|le couple}} $\lambda$.

% B08 / margin note / left beside example
% slanted
\margin{Digression}

% B09 / text / middle, underlined example and ring
\uline{Exemple} \deleted{\unsure{1)|)}} Soient $\Lambda$ un anneau \inserted{commut.}, $(M,\deleted{M^{\unsure{*|\vee}}}M^\vee,)$
\deleted{\unsure{$x\in M\times M^\vee$|$M\in X\times X^\vee$}} un couple de $\Lambda$-modules
libres de rang 2, accouplés par une forme bilinéaire
identifiant chacun d'eux au dual de l'autre. On
considère les \deleted{\unsure{relat.|relets|relat}} ensembles

% B10 / display / middle, (32), three lines
\[
 \left\{\begin{aligned}
 &\deleted{\unsure{e\in S=|s\in S=}}M^*\subset M\\
 &\deleted{\unsure{S^\vee=|s^\vee\in}}M^{\vee *}\subset M^\vee\\
 &S=M^*/\pm1,\quad S^\vee=M^\vee/\pm1
 \end{aligned}\right.\tag*{(32)}
\]

% B11 / text / middle right, parenthesis at (32)
(ens. des $x\in M$ \deleted{tel que} (resp. $\check{x}\in M^\vee$)
faisant partie d'une base)

% B12 / text / middle, introducing relation
avec la relation

% B13 / display / middle to lower, (33)
\[
 \begin{aligned}
 & R\subset S\times S^\vee\\
 & \rotatebox{90}{$=$}\\[-3pt]
 &\bigl\{(s,s^\vee)\mid\deleted{\unsure{\langle x,\check{x}\rangle\in\{\pm1\}|\langle x,\check{x}\rangle\in\pm1}}
 \ \exists(x,x^\vee)\in M^*\times M^{\vee *}\ \text{t.q.}\ 
 \underset{s=\dot x,\ s^\vee=\dot x^\vee.\ \deleted{\unsure{(x,\check{x})|(x,x^\vee)}}}
 {\langle x,\check{x}\rangle=1}\bigr\}
 \end{aligned}\tag*{(33)}
\]

% B14 / text / lower, struck paragraph opening
\deleted{\unsure{Mais|Ici} Nous nous intéressons aux applications}

% B15 / display / lower, struck bijection
\[
 \deleted{\varphi_0:S\simeq S^\vee}
\]

% B16 / text / lower, within cancelled block
\deleted{induites par des applications linéaires (notées par la \unsure{=|même} lettre)}

% B17 / display / lower, (34), within cancelled block
\[
 \deleted{\varphi_0:M\xrightarrow{\sim}M^\vee\qquad\Lambda\text{-linéaire}}\tag*{\deleted{(34)}}
\]

% B18 / text / lower, forms, within cancelled block
\deleted{qui correspondent donc aux formes bilinéaires non dégénérées}

% B19 / display / lower, pairing, within cancelled block
\[
 \deleted{\psi_0(x,y)=\deleted{\unsure{x\varphi|x}}(x,\varphi_0(y))\qquad M\times M\longrightarrow\Lambda\quad,}
\]

% B20 / text / lower, alternation, within cancelled block
\deleted{\deleted{satisfaisant la condition} qui sont \uline{alternées}}

% B21 / display / bottom, within cancelled block
\[
 \deleted{\psi_0(x,x)=0\qquad\text{i.e.}\quad (x,\varphi_0(x))=0}
\]

% B22 / text / bottom, within cancelled block
\deleted{et qui sont donc en corr. 1--1. avec les $e^\vee\in\bigwedge^2M^\vee=\det(M)^\vee$,}

% archive sheet number: 203
